% صفحة العنوان
\begin{titlepage}
\centering
\vspace*{1.5cm}
\begin{tcolorbox}[enhanced,colback=parchment,colframe=gold,boxrule=1.5pt,arc=4mm,
  borderline={0.5pt}{4pt}{gold},width=0.9\linewidth,center,halign=center,top=18pt,bottom=18pt]
  {\Large\color{headcolor} الشرح الميسَّر لـ}\par\medskip
  {\fontsize{34}{48}\selectfont\bfseries\color{maroon} حِرْزِ الأَمَانِي وَوَجْهِ التَّهَانِي}\par\medskip
  {\LARGE\color{headcolor} في القراءات السبع}\par\medskip
  {\Large (متن الشاطبية)}\par\bigskip
  {\large للإمام القاسم بن فِيرُّه الشاطبي (ت ٥٩٠ هـ)}
\end{tcolorbox}

\vspace{1.2cm}
{\large شرحٌ ميسَّر بيتًا بيتًا، مع بيان القراء بالألوان، والأمثلة القرآنية، والتحريرات}\par
\vspace{0.8cm}
{\large\color{maroon}\textbf{الجزء الأول}}\par
{\large مقدمة الناظم — باب الاستعاذة — باب البسملة — سورة أم القرآن}\par
{\large (الأبيات ١ – ١١٥)}\par

\vfill
{\normalsize مستخلَص من: الوافي في شرح الشاطبية لعبد الفتاح القاضي، وإتحاف البرية بتحريرات الشاطبية لحسن خلف الحسيني،\\ وتقريب الشاطبية، مع متن الشاطبية المضبوط}\par
\vspace{0.6cm}
\end{titlepage}
